% Part: lambda-calculus
% Chapter: introduction
% Section: computable-lambda

\documentclass[../../../include/open-logic-section]{subfiles}

\begin{document}

\olfileid[pt]{lam}{int}{lrp}
\olsection{Funções computáveis são \usetoken{S}{lambda definable}}

\begin{thm}
\ollabel{thm:computable-lambda}
Toda função parcial computável é !!{lambda definable}.
\end{thm}

\begin{proof}
Precisamos mostrar que toda função parcial computável~$f$ é
!!{lambda defined} por um termo lambda~$F$. Pelo teorema da forma normal de
Kleene, basta mostrar que toda função recursiva primitiva é !!{lambda
defined} por um termo lambda e, depois, que as funções !!{lambda definable}
são fechadas sob composições adequadas e busca ilimitada. Para mostrar que
toda função recursiva primitiva é !!{lambda defined} por um termo lambda,
basta mostrar que as funções iniciais são !!{lambda definable} e que as
funções parciais !!{lambda definable} são fechadas sob composição, recursão
primitiva e busca ilimitada.
\end{proof}

Usaremos uma notação mais convencional para tornar o restante da prova mais
legível. Por exemplo, escreveremos $M(x, y, z)$ em vez de $Mxyz$. Embora
isso seja sugestivo, lembre-se de que os termos do cálculo lambda não tipado
não têm aridades associadas; assim, para o mesmo termo~$M$, faz tanto sentido
escrever $M(x,y)$ quanto $M(x,y,z,w)$. Mas o uso dessa notação indica que
estamos tratando $M$ como função de três variáveis e ajuda a tornar mais
claras as intenções por trás das definições. De modo semelhante, diremos
``defina $M$ por $M(x,y,z) = \dots$'' em vez de ``defina $M$ por $M =
\lambd[x][\lambd[y][\lambd[z][\dots]]]$.''

\end{document}

